\section{Isolating a compact piece}
\label{sec:isolation}
A bounded component need not be separated from the rest of a closed set.
The following lemma enlarges it to a separated compact piece without
losing connectedness of the complement after removal.

\begin{lemma}[Isolating a compact piece]
\label{lem:compact-isolation}
Let $H\subset\mathbb R^2$ be closed, locally of finite length, and equal to the
support of $\mathcal H^1\!\llcorner H$.  Every nonempty bounded connected
component of $H$ is contained in a nonempty compact subset $A\subset H$ such
that, with $B=H\setminus A$,
\[
 B\text{ is closed},\qquad
 \operatorname{dist}(A,B)>0,\qquad
 0<\mathcal H^1(A)<\infty.
\]
Here the distance to the empty set is interpreted as infinity.  If
$\mathbb R^2\setminus H$ is connected, $\mathbb R^2\setminus B$ is connected.
\end{lemma}

\begin{proof}
Let $C$ be the given component.  It is compact, so choose $R$ with
$C\subset B_R$ and set $X=H\cap\overline{B_R}$.  The component in $X$ of any
$p\in C$ is exactly $C$.

We use the compact-space fact that a component is the intersection of all its
clopen neighborhoods.  For completeness, in a compact metric space let $Q_n$
be the points joined to $p$ by a finite chain with consecutive distances less
than $1/n$.  Each $Q_n$ is clopen.  Their nested intersection $Q$ is connected:
otherwise the two compact pieces of a separation admit disjoint neighborhoods
at positive distance; for large $n$, compactness puts $Q_n$ in their union,
where a $1/n$-chain cannot join the two pieces.  Thus $Q$ is the component of
$p$, proving the stated fact.

For each $z\in X\cap\partial B_R$, choose a clopen neighborhood of $C$ in
$X$ which omits $z$.  Finitely many of their complements cover
$X\cap\partial B_R$; take $A$ to be the intersection of the corresponding
neighborhoods.  If $X\cap\partial B_R$ is empty, take $A=X$.  In either case
$A$ is compact, contains $C$, and is disjoint from $\partial B_R$.  Its positive
distance from that circle shows that its relative openness in $X$ also gives
relative openness in $H$.  Consequently $B=H\setminus A$ is closed, and its
distance from compact $A$ is positive.  Local finite length gives
$\mathcal H^1(A)<\infty$; the support assumption and relative openness give
$\mathcal H^1(A)>0$.

Finally, $D=\mathbb R^2\setminus H$ is dense because $H$ has locally finite
length.  A set lying between a connected set and its closure is connected.
The larger complement $\mathbb R^2\setminus B$ lies between $D$ and
$\overline D=\mathbb R^2$, which proves the last assertion.
\end{proof}


The compact piece $A$ may contain several connected components of $H$.
What matters in Section~\ref{sec:compact-piece} is its positive separation
from the fixed remainder $B$, which permits every sufficiently small
translation of the whole piece.
